\subsection{آزمایش‌های ثانویه و اکتشافی مدل‌های عصبی}

سه آزمایش این زیربخش ثانویه و اکتشافی‌اند و روی همان آزمون تقسیم استاندارد قفل‌شده اجرا شده‌اند، پس بار تأییدی ندارند. انتخاب مدل در آن‌ها هم فقط با اعتبارسنجی انجام شد. جدول زیر معیارهای آزمون کانونی چهار مدل دو آزمایش نخست را کنار هم می‌گذارد.

\begin{table}[htbp]
\centering
\renewcommand{\arraystretch}{1.35}
\small
\caption{آزمایش‌های ثانویه و اکتشافی معماری که از مدل‌های الحاقی تثبیت‌شده بهتر نشدند}
\label{tab:negative}
\setlength{\tabcolsep}{3pt}
\begin{tabular}{>{\raggedleft\arraybackslash}m{3.6cm}rrrrr>{\centering\arraybackslash}m{4.2cm}}
\hline
مدل & \lr{AUC} & دقت متعادل & \lr{F1} ماکرو & \lr{Brier} & \lr{ECE} & توضیح \\
\hline
ترکیب مشترک یادگیرنده (\lr{H}) & 0.9258 & 0.6574 & 0.6840 & 0.2759 & 0.1159 & دوره‌ی انتخاب‌شده 4، مرحله‌ی دوم \\
\hline
ترکیب مشترک + نشانگر (\lr{I}) & 0.9162 & 0.6730 & 0.6978 & 0.2786 & 0.1152 & دوره‌ی انتخاب‌شده 6، مرحله‌ی دوم \\
\hline
کنترل ترکیب تثبیت‌شده (\lr{J0}) & 0.9298 & 0.7954 & 0.6852 & 0.3770 & 0.1676 & دوره‌ی انتخاب‌شده 1 \\
\hline
شرطی \lr{FiLM} (\lr{J1}) & 0.9276 & 0.7897 & 0.6885 & 0.3614 & 0.1536 & دوره‌ی انتخاب‌شده 1 \\
\hline
\multicolumn{7}{p{0.96\textwidth}}{\footnotesize هیچ‌یک از این مدل‌ها نتیجه‌ی اصلی نیست. دقت متعادل بالاتر در \lr{J0} و \lr{J1} ناشی از پیش‌بینی تهاجمی کلاس‌های کم‌فراوان تحت تابع زیان وزن‌دار است و با \lr{AUC}، \lr{F1} ماکرو و کالیبراسیون بدتر همراه است؛ بنابراین نباید به‌عنوان برتری خوانده شود.} \\
\hline
\end{tabular}
\end{table}

\subsubsection{ترکیب مشترک یادگیرنده (\lr{Task~9})}

در این آزمایش دو رمزگذار تثبیت‌شده‌ی تصویری و عروقی به بردارهای فشرده نگاشت شدند، یک شاخه‌ی تعامل روی آن دو ساخته شد و یک سر چندلایه تصمیم گرفت (مدل \lr{H})؛ مدل \lr{I} همان ساختار را با یک شاخه‌ی نشانگری کوچک تکرار کرد. هیچ‌یک از این دو به ترکیب الحاقی تصویر و عروق نمی‌رسد. کاهش \lr{H} نسبت به آن ترکیب پشتیبانی نمی‌شود و افزودن شاخه‌ی نشانگری به \lr{H} کاهش پشتیبانی‌شده‌ای در رتبه‌بندی می‌دهد (جدول شکل مربوط).

معیارهای دیگر همان مقایسه‌ها یک‌سو نیستند. مدل \lr{H} در برابر ترکیب الحاقی در دقت متعادل و \lr{F1} ماکرو و کیفیت احتمال عقب می‌ماند، در حالی که شاخه‌ی نشانگری در \lr{I} هیچ‌یک از این معیارها را معنادار جابه‌جا نمی‌کند؛ یعنی آنچه در \lr{I} افت کرد رتبه‌بندی بود، نه کالیبراسیون. منحنی‌های یادگیری نشان می‌دهند پس از آزادسازی وزن‌های رمزگذار، خطای آموزش یکنوا پایین رفت و خطای اعتبارسنجی بالا رفت، و کمینه‌ی خطای اعتبارسنجی هر دو مدل پیش از آزادسازی بود. آنچه این آزمایش نشان می‌دهد بیش‌برازش سر مشترک روی نمونه‌های آموزش است، نه برتری معماری ترکیبی. این آزمایش ثانویه و اکتشافی است.

\begin{figure}[htbp]
  \centering
  \imgbox[1.0\textwidth]{figures/report/appendix/N1a_task9_loss.pdf}
  \caption{خطای آموزش و اعتبارسنجی مدل‌های تسک نه به‌ازای هر دوره؛ رمزگذارها پس از دوره‌ی سوم آزاد شده‌اند.}
  \label{fig:U6}
\end{figure}

\begin{figure}[htbp]
  \centering
  \imgbox[1.0\textwidth]{figures/report/appendix/N1b_task9_validation.pdf}
  \caption{توان تمایز مدل‌های تسک نه روی اعتبارسنجی به‌ازای هر دوره.}
  \label{fig:U7}
\end{figure}

\begin{figure}[htbp]
  \centering
  \imgbox[1.0\textwidth]{figures/report/appendix/N2a_task10_loss.pdf}
  \caption{خطای آموزش و اعتبارسنجی مدل‌های تسک ده به‌ازای هر دوره.}
  \label{fig:U8}
\end{figure}

\begin{figure}[htbp]
  \centering
  \imgbox[1.0\textwidth]{figures/report/appendix/N2b_task10_validation.pdf}
  \caption{توان تمایز مدل‌های تسک ده روی اعتبارسنجی به‌ازای هر دوره.}
  \label{fig:U9}
\end{figure}

\begin{figure}[htbp]
  \centering
  \imgbox[1.0\textwidth]{figures/report/appendix/B2a_film_modulation.pdf}
  \caption{تعدیل یادگرفته‌شده‌ی \lr{FiLM} روی آزمون؛ خط‌چین‌ها حالت همانی را نشان می‌دهند.}
  \label{fig:U10}
\end{figure}

\begin{figure}[htbp]
  \centering
  \imgbox[1.0\textwidth]{figures/report/appendix/B2b_biomarker_sensitivity.pdf}
  \caption{تغییر مطلق احتمال پیش‌بینی مدل تثبیت‌شده هنگام جایگزینی هر نشانگر با میانگین آموزشی خود.}
  \label{fig:U11}
\end{figure}

\begin{figure}[htbp]
  \centering
  \imgbox[1.0\textwidth]{figures/report/main/R4_multiclass_auc.pdf}
  \caption{سطح زیر منحنی چندکلاسه‌ی مدل‌های تسک سیزده در برابر مدل‌های مرجع، روی آزمون کانونی (1331 تصویر).}
  \label{fig:U12}
\end{figure}

\begin{figure}[htbp]
  \centering
  \imgbox[1.0\textwidth]{figures/report/main/R4_balanced_accuracy.pdf}
  \caption{دقت متعادل مدل‌های تسک سیزده در برابر مدل‌های مرجع، روی آزمون کانونی (1331 تصویر).}
  \label{fig:U13}
\end{figure}

\begin{figure}[htbp]
  \centering
  \imgbox[1.0\textwidth]{figures/report/main/R4_macro_f1.pdf}
  \caption{اف‌یک ماکروی مدل‌های تسک سیزده در برابر مدل‌های مرجع، روی آزمون کانونی (1331 تصویر).}
  \label{fig:U14}
\end{figure}

\begin{figure}[htbp]
  \centering
  \imgbox[1.0\textwidth]{figures/report/main/R4_brier.pdf}
  \caption{امتیاز \lr{Brier} مدل‌های تسک سیزده در برابر مدل‌های مرجع، روی آزمون کانونی (1331 تصویر).}
  \label{fig:U15}
\end{figure}

\begin{figure}[htbp]
  \centering
  \imgbox[1.0\textwidth]{figures/report/main/R4_ece.pdf}
  \caption{خطای کالیبراسیون انتظاری مدل‌های تسک سیزده در برابر مدل‌های مرجع، روی آزمون کانونی (1331 تصویر).}
  \label{fig:U16}
\end{figure}

\begin{figure}[htbp]
  \centering
  \imgbox[1.0\textwidth]{figures/report/main/R5_task13_paired_auc.pdf}
  \caption{چهار مقایسه‌ی زوجی تسک سیزده، با بازه‌ی اطمینان از بوت‌استرپ طبقه‌محور با ده هزار تکرار و بذر ثابت روی آزمون کانونی (1331 تصویر).}
  \label{fig:U17}
\end{figure}

\begin{figure}[htbp]
  \centering
  \imgbox[1.0\textwidth]{figures/report/appendix/N3a_task13_loss.pdf}
  \caption{خطای آموزش و اعتبارسنجی مدل تسک سیزده به‌ازای هر دوره؛ ستون‌های فقرات پس از دوره‌ی سوم آزاد شده‌اند.}
  \label{fig:U18}
\end{figure}

\begin{figure}[htbp]
  \centering
  \imgbox[1.0\textwidth]{figures/report/appendix/N3b_task13_validation.pdf}
  \caption{توان تمایز مدل تسک سیزده روی اعتبارسنجی به‌ازای هر دوره.}
  
\end{figure}


\subsubsection{شرطی‌سازی نشانگری با تعدیل کران‌دار (\lr{Task~10})}

مدل \lr{J0} کنترل این آزمایش است؛ دو بازنمایی تثبیت‌شده به یک سر مشترک داده می‌شوند و هیچ نشانگری وارد مسیر نمی‌شود. مدل \lr{J1} همان ساختار را با تعدیل \lr{FiLM} روی شاخه‌ی عروقی تکرار می‌کند؛ پارامترها باقی‌مانده و کران‌دار نوشته شدند تا مقدار اولیه نزدیک همانی بماند. تفاضل رتبه‌بندی \lr{J1} و \lr{J0} بازه‌ی شامل صفر دارد، ولی کیفیت احتمال در \lr{J1} بهتر است و دقت متعادل و \lr{F1} ماکرو بی‌تفاوت می‌مانند.

این بهبود کالیبراسیون به نشانگرها نسبت داده نمی‌شود، چون تشخیص‌های همین آزمایش خلاف آن را نشان می‌دهند. تعدیل یادگرفته‌شده روی آزمون تقریباً همانی ماند و آزمون حساسیت تک‌نشانگری هم همین را تأیید می‌کند؛ جانشین‌کردن هر نشانگر با میانگین آموزشی‌اش احتمال‌های تثبیت‌شده‌ی \lr{J1} را تغییر محسوسی نمی‌داد و شمار تغییر تصمیم در هر پنج نشانگر صفر بود. توضیح محتمل آن بهبود یک محدودیت طرح تطبیقی است، نه یک یافته؛ ساخت ماژول‌های شرطی جریان تصادفی را جابه‌جا می‌کند، پس مقدار اولیه‌ی سر ترکیب در \lr{J0} و \lr{J1} یکسان نیست و این دو مدل یک مقایسه‌ی کاملاً زوج‌شده نیستند. در منحنی‌های یادگیری هم هر دو مدل دوره‌ی نخست را انتخاب کردند و در پایان خطای آموزش پایین آمد در حالی که خطای اعتبارسنجی بالا رفت. این آزمایش هم ثانویه و اکتشافی است.



\subsubsection{استواری معماری نتیجه‌ی عروقی (\lr{Task~13})}

این آزمایش یک ستون فقرات تصویری دیگر می‌آزماید؛ شبکه‌ی توجه نرم دو‌ستون‌فقرات با \lr{ResNet50} و \lr{EfficientNet-B4} که یک بازنمایی توجهی می‌سازد (مدل \lr{L})، همان بازنمایی با طبقه‌بند جدولی (مدل \lr{M0}) و همان بازنمایی به‌همراه نمایش عروقی (مدل \lr{M1}). طبقه‌بند خطی سرتاسری \lr{L} از طبقه‌بند جدولی روی همان بازنمایی ضعیف‌تر است و این دو هم‌خانواده نیستند، پس \lr{L} مستقیماً با مدل‌های بازنمایی قابل مقایسه نیست.

مقایسه‌ی کلیدی \lr{M0} در برابر بازنمایی تصویری اصلی است که کاهش پشتیبانی‌شده‌ای در رتبه‌بندی و \lr{F1} ماکرو و کیفیت احتمال می‌دهد؛ یعنی بازنمایی توجهی آزموده‌شده در این پروتکل از بازنمایی تک‌ستون‌فقرات اصلی بدتر بود. این نتیجه یک آبلاسیون کنترل‌شده‌ی معماری نیست، چون دو خط لوله در یادگیرنده و رژیم آموزش هم تفاوت دارند. افزودن عروق در همین خانواده نتیجه‌ی جهت‌داری می‌دهد و بهبود پشتیبانی‌شده‌ای در رتبه‌بندی و \lr{F1} ماکرو و کیفیت احتمال می‌سازد، ولی همین \lr{M1} در برابر مدل‌های اصلی تفاوت معناداری با آن‌ها ندارد. پس مسیر عروقی در معماری دوم هم مکمل می‌ماند، بدون شاهد برتری خودِ آن معماری.

دوره‌ی برگزیده‌ی \lr{L} پیش از بازکردن وزن‌های ستون فقرات قرار داشت و کمینه‌ی خطای اعتبارسنجی در دوره‌های بعد آمد؛ یعنی مرحله‌های تنظیم دقیق هرگز به سود نرسیدند و مدل برگزیده عملاً یک مدل مرحله‌ی نخست است. این آزمایش ثانویه و اکتشافی است.





\subsubsection{جمع‌بندی بیش‌برازش}

در هر سه آزمایش نورونی این زیربخش، خطای آموزش کاهش یافت در حالی که عملکرد اعتبارسنجی افت کرد و کمینه‌ی خطای اعتبارسنجی پیش از بهترین دوره‌ی تنظیم دقیق یا در همان دوره‌ی نخست بود. سرهای عصبی این آزمایش‌ها روی نمونه‌های آموزش برای ظرفیت‌شان داده‌ی کافی ندارند و همین نتیجه‌ی منفی هر سه را توضیح می‌دهد. این سه آزمایش ثانویه و اکتشافی‌اند و هیچ‌کدام جای نتیجه‌ی مقایسه‌ی اصلی را نمی‌گیرد.
